\documentclass[10pt,a4paper]{amsart}
\usepackage[margin=19mm]{geometry}
\usepackage{amsmath,amssymb,mathtools}
\usepackage[scr=rsfs,cal=euler]{mathalfa}
\usepackage{xfrac}
\usepackage[unicode,hidelinks,bookmarks=false]{hyperref}
\newcommand{\ie}{i.e.\ }
\newcommand{\cf}{cf.\ }
\newcommand{\resp}{respectively\ }
\newcommand{\Subsection}{\subsection}
\providecommand{\nobreakdash}{\nobreak\hbox{-}}
\setlength{\emergencystretch}{2em}
\multlinegap0pt
\usepackage{enumerate}
\usepackage{float}
\usepackage{mathtools}
\usepackage{tcolorbox}
\usepackage{tikz}
\newtheorem{lemma}{Lemma}
\DeclareMathOperator{\rank}{rank}
\title{Bootstrapping bilinear relations of discrete Painlev\'e systems from 5d gauge theories}
\author{Minsung Kim, Xin Wang}
\date{Source: arXiv:2608.15756v1 (CC BY 4.0)}
\begin{document}
\maketitle

\noindent\emph{Source and licence.} Bootstrapping bilinear relations of discrete Painlev\'e systems from 5d gauge theories. Version arXiv:2608.15756v1 (CC BY 4.0), distributed by arXiv under the Creative Commons Attribution 4.0 International licence, http://creativecommons.org/licenses/by/4.0/, as recorded in the arXiv licence metadata for this exact version and shown on the abstract page https://arxiv.org/abs/2608.15756v1. Full licence notice: \texttt{licenses/arxiv-2608.15756-CC-BY-4.0.txt}.\par\smallskip

\noindent\textit{Editorial scope.} Counted unit: the article's lemma lemma (source0.tex lines 1477--1484). The article's complete original proof of it is delivered with it (source0.tex lines 1485--1516, one \texttt{proof} environment of 32 lines). No mathematics is written, repaired or re-derived: the statement and the proof are the article's own lines, in the article's own notation, and no retained line is substituted or re-flowed.  No further source-local context is needed: every cross-reference of every retained line resolves inside the two delivered spans.  Nothing was omitted from any retained span, so the package carries no omission record.  The retained lines cite 1 works, whose entries are transcribed verbatim from the article's own bibliography source in the e-print and delivered with short editorial labels recorded in the item's reference mapping.  No retained line contains a figure environment, an image-inclusion command or drawing code, so no image is delivered and the package declares no asset dependency.  Editorial formatting only, in addition: an AMS \texttt{amsart} preamble that restates the article's own declarations and macros used by the retained lines, namely the environments \texttt{lemma} on separate counters, together with the standard packages those lines need.  The retained environments print in this document as Lemma 1 in order of appearance, whereas the article prints the same items with the numbering of its own full document.  The complete native source of the article as distributed in the arXiv e-print for this exact version is bundled unchanged as \texttt{sources/207781/source0.tex}.
\begin{lemma}
    Let $ \{\alpha_{ia}\}_{i=1,2,3}^{a=1,2} $ and $ \{\beta_{ia}\}_{i=1,2,3}^{a=1,2} $ be two sets of nonzero $ E_8 $ root vectors satisfying the constant-sum conditions
    \begin{align}
        \alpha_{i1}+\alpha_{i2}=\lambda \, , \quad
        \beta_{i1}+\beta_{i2}=\sigma \, ,
    \end{align}
    for all $ i $, and set $ x_i = \alpha_{i1}-\alpha_{i2} $ and $ y_i=\beta_{i1}-\beta_{i2} $. If the Gram matrices of $ x_i $ and $ y_i $ agree, $ \langle x_i, x_j \rangle = \langle y_i, y_j \rangle $ for all $ i,j $, then there exists $ w \in W(E_8) $ such that $ w(\alpha_{ia}) = \beta_{ia} $.
\end{lemma}

\begin{proof}
    Since $ \alpha_{ia} $ and $ \beta_{ia} $ are nonzero $ E_8 $ roots, $ \lVert \alpha_{ia} \rVert^2 = \lVert \beta_{ia} \rVert^2 = 2 $, and  $ \langle \lambda, x_i \rangle = \langle \sigma, y_i \rangle = 0 $. The roots are decomposed as 
    \begin{align}\label{eq:E8root-decompose}
        \alpha_{i1} = \frac{\lambda+x_i}{2} \, , \quad
        \alpha_{i2} = \frac{\lambda-x_i}{2} \, , \quad
        \beta_{i1} = \frac{\sigma+y_i}{2} \, , \quad
        \beta_{i2} = \frac{\sigma-y_i}{2} \, .
    \end{align}
    Squaring \eqref{eq:E8root-decompose} gives $ \lVert \lambda \rVert^2 + \lVert x_i \rVert^2 = \lVert\sigma\rVert^2 + \lVert y_i \rVert^2 = 8 $. Combined with $ \lVert x_i \rVert^2 = \lVert y_i \rVert^2 $, this yields $ \lVert \lambda \rVert^2 = \lVert \sigma \rVert^2 $. Then a direct computation gives 
    \begin{align}\label{eq:app-inner}
       \langle \alpha_{ia}, \alpha_{jb} \rangle = \frac{\lVert \lambda \rVert^2 \pm \langle x_i, x_j \rangle}{4} = \frac{\lVert \sigma \rVert^2 \pm \langle y_i, y_j \rangle}{4} = \langle \beta_{ia}, \beta_{jb} \rangle \, ,
    \end{align}
    where the sign is plus if $ a=b $ and minus if $ a\neq b $. 

    Let $ \Sigma $ be the $ E_8 $ root system and $ \Xi_A, \Xi_B \subset \Sigma $ be the smallest root subsystems generated by $ \{\alpha_{ia}\} $ and $ \{\beta_{ia}\} $, respectively. Since $ \alpha_{ia} \in \operatorname{Span}\{\lambda, x_1, x_2, x_3 \} $ and $ \beta_{ia} \in \operatorname{Span}\{\sigma, y_1,y_2,y_3\} $, both root subsystems have rank at most four. By \eqref{eq:app-inner}, the assignment $ f(\alpha_{ia}) = \beta_{ia} $ extends to an isomorphism $ f : \Xi_A \to \Xi_B $ of root systems preserving the inner product. By the classification of root subsystems of $ E_8 $ \cite{Oshima:2007}, possible subsystems of rank at most four are
    \begin{align}
        A_1,\ A_2,\ A_3,\ A_4,\ D_4,\ 2A_1,\ 3A_1,\ 4A_1,\ A_2+A_1,\ A_2+2A_1,\ 2A_2,\ A_3+A_1 .
    \end{align}
    Let $ \iota_A : \Xi_A \hookrightarrow \Sigma $ and $ \iota_B : \Xi_B \hookrightarrow \Sigma $ be the inclusion maps. For every type on the list except $ 4A_1 $, all embeddings of that type into $ \Sigma $ lie in a single $ W(E_8) $-orbit \cite{Oshima:2007}. Consequently, provided $ \Xi_A \not\cong 4A_1 $, the two embeddings $ \iota_A $ and $ \iota_B \circ f $ of $ \Xi_A $ into $ \Sigma $ are $ W(E_8) $-related, i.e., there exists $ w \in W(E_8) $ with $ \iota_B \circ f = w \circ \iota_A $. Evaluating on the generators yields
    \begin{align}
        w(\alpha_{ia}) = \beta_{ia}
    \end{align}
    for all $ i $ and $ a $.

    It remains to exclude the case $ \Xi_A \cong 4A_1 $. Assuming this for
    contradiction, write
    \begin{align}
        \Xi_A = \{ \pm \gamma_1, \pm \gamma_2, \pm \gamma_3, \pm \gamma_4 \} \, ,
        \qquad \langle \gamma_j, \gamma_k \rangle = 2\delta_{jk} \, .
    \end{align}
    Each $ \alpha_{ia} $ is then of the form $ \pm \gamma_k $ for some $ k $, and the constant-sum condition $ \alpha_{i1}+\alpha_{i2}=\lambda $ leaves three possibilities for $ \lambda $: $ \lambda = 0 $, $ \pm 2\gamma_k $, or $ \pm\gamma_j \pm \gamma_k $ with $ j\neq k $. If $ \lambda = 0 $, then $ \alpha_{i2} = -\alpha_{i1} $ for each $ i $, so the six roots $ \{\alpha_{ia}\} $ span at most three dimensions, contradicting $ \rank\Xi_A = 4 $. If $ \lambda \neq 0 $, the pair $ \{\alpha_{i1}, \alpha_{i2}\} $ is uniquely determined by $ \lambda $, so the six vectors span at most two dimensions, again contradicting $ \rank\Xi_A = 4 $. Either way, $ \Xi_A \not\cong 4A_1 $, completing the proof.
\end{proof}

\begin{thebibliography}{99}
\bibitem[Oshima]{Oshima:2007} T.~Oshima, \emph{{A classification of subsystems of a root system}}, \href{https://arxiv.org/abs/math/0611904}{{\ttfamily math/0611904}}.
\end{thebibliography}
\end{document}
